% Lösungen Einheit 1 von 5 – Bruch als Anteil (zusammenbau v0.1)
\einheitenkopf{Einheit 1 von 5 · Bruch als Anteil}

\erg{10}{a) ja \quad b) $\frac{3}{5}$ \quad c) $5$ von $6$ Teilen \quad d) $\frac{5}{8}$ (in Achteln: A zwei, B zwei, D eins) \quad e) Q \quad f) $\frac{4}{9}$ \quad g) $\frac{3}{9}$}

10c)

\bruchrechteck{5}{6}

\erg{11}{a) alle $24$ Bonbons der Tüte \quad b) $8$ \quad c) $36 : 4 = 9$, $9 \cdot 3 = 27$ \quad d) $27$ € \quad e) $12 : 3 \cdot 2 = 8$ Stücke gegessen, $12 - 8 = 4$ Stücke übrig \quad f) $1{,}5$ kg $= 1\,500$ g, $1\,500 : 5 \cdot 2 = 600$ g \quad g) $120 : 4 \cdot 3 = 90$ l sind drin, es fehlen $120 - 90 = 30$ l}
\erg{12}{a) $8 \cdot 3 = 24$ Kekse}
\erg{13}{a) Er hat den Zähler als Anzahl der Kästchen genommen. Richtig: $12 : 3 \cdot 2 = 8$ Kästchen. \quad b) Bei $4$ Kindern: Je mehr Kinder teilen, desto mehr Teile gibt es, und jedes Teil wird kleiner. \quad c) $\frac{5}{9}$ \quad d) $90 : 3 \cdot 2 = 60$ min gespielt, $90 - 60 = 30$ min auf der Bank}
